\begin{exsol}{15}{\MR{} Food for Thought 4.3}{119}{Exam}{A3}
\asked Explain why consumer-electronics leaders (Sony, then Matsushita) bought record, film and
television producers around 1987--1990.

\method Complementarity as rising marginal profitability, $\partial^2\pi/\partial x_i\partial
x_j\ge0$ (\MR{} Ch.~4, p.~108 and n.~8). Strong complementarity and design attributes (p.~109).
Multiple coherent patterns and coordination failure (pp.~111--113). Design decisions:
synchronization and the communication of design variables such as a technical standard
(pp.~104--106; the HDTV standard, p.~108). Economies of scope and core competencies (p.~107), as the
contrasting motive.

\solution
\subparagraph{The main motive: complementarity between hardware and content.} Let $x_h$ be the
firm's effort in a playback format (players, recorders, televisions) and $x_c$ the content
available in that format (records, films, programmes). Consumers buy a player for what they can
play on it, and content sells only to people who own players. So more content raises the marginal
profit of selling hardware, and more hardware raises the marginal profit of releasing content:
$\partial^2\pi/\partial x_h\partial x_c\ge0$, the definition of p.~108. A firm doing only one side
captures only part of the joint return. Owning both lets it internalize the externality: release
content in its new format early, price hardware to build the installed base, and accept losses on
one side for gains on the other.

\subparagraph{Design decisions and synchronization.} A new format is a design variable with both
attributes of p.~91. Its success requires a predictable fit --- players, discs or tapes and a
content catalogue in the same standard, available at the same time --- and a misfit (a format
without titles at launch) is very costly, while small differences in technical quality matter
much less. Like the car launch date (p.~104), the launch of a format is best coordinated by
announcing the design variables (the standard, the date) to all parties. The text's own example is
HDTV: Japan's push was grounded in electronics and displays, and a central standard focused the
effort (p.~108). Owning a studio and a record label puts both sides of the design decision under
one management.

\subparagraph{Coordination failure and multiple coherent outcomes.} With strong complementarities
there are several self-reinforcing outcomes (pp.~111--113): ``no content, so no players, so no
content'' is coherent, and so is ``rich catalogue, large installed base''. Independent studios,
each waiting for others, can leave a new format in the low outcome or choose a rival standard. The
lesson the electronics firms drew from home-video format contests --- where the format with more
available titles prevailed --- is that hardware quality alone does not decide the outcome. Owning
a content library is a way to make the jump to the good outcome rather than wait for it.

\subparagraph{What the motive is not.} It is not mainly economies of scope or shared core
competencies (p.~107). Making films and records draws on creative, talent-based capabilities quite
unlike precision electronics manufacturing; little is shared on the production side. The
complementarity is in \emph{demand}, not in production costs. Other motives may also have
mattered (diversification, growth), but they do not explain why these particular industries were
chosen.

\subparagraph{The trade-off.} Complementarity between activities does not by itself require
common ownership: licensing and contracts with independent labels and studios can coordinate a
launch, and many formats have succeeded that way. Ownership pays when coordination must be
repeated across formats and is hard to contract for. And it carries a cost that Chapter~4
predicts: a Hollywood studio and a Japanese electronics maker are each coherent systems with very
different incentives, careers and decision rights, and the parent's top management lacks
information about the creative business (innovation attributes, p.~92). Imposing one firm's
policies on the other breaks the internal fit of the acquired business.

\interp Complementarities explain the direction of expansion: firms grow into the activities that
raise the marginal return on what they already do. They also explain why such expansions are hard
to manage --- each business is itself a system of complementary practices, and combining two
coherent systems does not make a coherent whole.

\verify Comparative statics with complements (p.~108, n.~8): a shift that raises the return to
hardware (say, a new disc format) raises the optimal level of content, and vice versa. The
prediction is that acquisitions cluster around new formats --- consistent with the timing, as CDs,
home video and the prospect of HDTV were all live issues in 1987--1990. If hardware and content
were independent ($\partial^2\pi/\partial x_h\partial x_c=0$), the acquisitions would have no
strategic logic beyond financial diversification.

\answer{Because consumer electronics and recorded content are complements in demand: content
raises the return to players and players raise the return to content. Owning studios and labels
lets the firm coordinate the design decision of a new format (standard, launch timing, catalogue)
and escape the ``no content, no players'' coordination failure. The motive is complementarity, not
shared production competencies; the cost is managing a business whose coherent organization is
very different from the parent's.}
\end{exsol}
